% !PRINTED 324
\weberpage{324}
Any ternary form of degree $n$, $f(x_1, x_2, x_3)$, set equal to zero,
represents a curve of the $n$th order, which has $n$ points of intersection
with a straight line. We denote this curve briefly as the curve $f$.

By the substitution (1), $f$ passes into another form of the same degree in the variables $y_1, y_2, y_3$:

\weq{3}{f(x_1, x_2, x_3) = \varphi(y_1, y_2, y_3),}

\noindent and $\varphi = 0$ represents the same curve as $f = 0$, referred to
the coordinate system $y_1, y_2, y_3$.

If we form the derivatives of the identity (3) with respect to the
variables $y_1, y_2, y_3$, it follows, having regard to (1):

\weq{4}{\begin{aligned}
\varphi'(y_1) &= \alpha_1^{(1)} f'(x_1) + \alpha_2^{(1)} f'(x_2) + \alpha_3^{(1)} f'(x_3) \\
\varphi'(y_2) &= \alpha_1^{(2)} f'(x_1) + \alpha_2^{(2)} f'(x_2) + \alpha_3^{(2)} f'(x_3) \\
\varphi'(y_3) &= \alpha_1^{(3)} f'(x_1) + \alpha_2^{(3)} f'(x_2) + \alpha_3^{(3)} f'(x_3).
\end{aligned}}

\begin{center}
\S.\,86.

\gesperrt{Singular Points. Points of Inflection. Double Tangents.}
\end{center}

If for a point $x$ the three equations

\weq{1}{f'(x_1) = 0, \quad f'(x_2) = 0, \quad f'(x_3) = 0}

\noindent are satisfied simultaneously, then by \gesperrt{Euler's theorem}

\weq{2}{n f = x_1 f'(x_1) + x_2 f'(x_2) + x_3 f'(x_3)}

\noindent (Bd. I, \S.\,17) this point lies on the curve $f$. The equations \S.\,85,(4)
show that at the same point $\varphi'(y_1), \varphi'(y_2), \varphi'(y_3)$
must also vanish, hence that the property of a point characterised by the
equations (1) is preserved under linear transformation, and therefore, as we
also express it, belongs to the \gesperrt{invariant} properties.

Not on every curve of the $n$th order do points occur which satisfy
the conditions (1), as is shown, e.g., by the assumption $f(x_1, x_2, x_3)$
$= x_1^n + x_2^n + x_3^n$. Rather, if even a single such point is to exist,
an equation between the coefficients of $f$ must be satisfied, which one
obtains by elimination of $x_1, x_2, x_3$ from the three equations (1) (Bd. I, \S.\,50).
One can represent this relation in the form that a certain integral
rational and homogeneous function of all the coefficients of $f$ must be
equal to zero, and this function
%CONT>>
